% GENERATED FILE. Edit sanskrit/pronunciation-reference.md, then run npm run generate:books.

\chapter*{The Devanagari script and Sanskrit sounds --- a reference}
\addcontentsline{toc}{chapter}{Devanagari \& Sanskrit sounds (reference)}
\markboth{Devanagari}{Devanagari}

A \textbf{reference}, not a chapter — the lessons teach reading through real words. This page gathers the system and the points Sanskrit needs even from a Devanagari reader.

\section*{How Devanagari works (the three facts)}

\begin{enumerate}
\raggedright
\setlength{\itemsep}{0.35em}
  \item \textbf{Consonants carry a built-in \textquotedblleft{}a\textquotedblright{}.} \sk{क} = ka, \sk{न} = na, \sk{म} = ma — left to right under a connecting top line.
  \item \textbf{A vowel sign changes it:} \sk{का} kā · \sk{कि} ki · \sk{की} kī · \sk{कु} ku · \sk{के} ke · \sk{को} ko. A word-initial vowel has its own letter: \sk{अ} a · \sk{आ} ā · \sk{इ} i · \sk{उ} u · \sk{ए} e · \sk{ओ} o.
  \item \textbf{A \emph{halant} (\sk{्}) removes the vowel}, and the bare consonant joins the next as a conjunct: \sk{स्} + \sk{त} $\to$ \sk{स्त}, \sk{न्} + \sk{य} $\to$ \sk{न्य}.
\end{enumerate}

\section*{What Sanskrit needs even from a Devanagari reader}

\begin{itemize}
\raggedright
\setlength{\itemsep}{0.35em}
  \item \textbf{The visarga \sk{ः} (ḥ)} — a soft breathed echo of the preceding vowel at a word's end (\emph{namaskāraḥ}, \emph{rāmaḥ}), usually the masculine-singular ending, so grammatical as well as phonetic.
  \item \textbf{The anusvāra \sk{ं} (ṃ, also written ṅ)} — a dot above the line, nasalising the vowel or standing for a nasal before another consonant; very frequent in Sanskrit.
  \item \textbf{The vocalic ṛ \sk{ऋ}} — Sanskrit treats ṛ as a \emph{vowel}: \emph{mātṛ} \textquotedblleft{}mother,\textquotedblright{} \emph{pitṛ} \textquotedblleft{}father\textquotedblright{} (the words that proved Sanskrit's kinship with Latin \emph{māter}/\emph{pater} and English \emph{mother}/\emph{father}), \emph{Ṛgveda}.
  \item \textbf{Sandhi} — fixed sound-rules fuse words at their seams: \emph{su} + \emph{āgata} $\to$ \emph{svāgata}. The lessons flag each fusion where it happens.
\end{itemize}

\section*{Sanskrit in the family}

Sanskrit is the classical language of India — the tongue of the Vedas, of Kālidāsa, of Hindu, Buddhist, and Jain learning — and the ancestor of the modern Indo-Aryan languages (Hindi, Marathi, Punjabi, Bengali, and more). It is a senior branch of \textbf{Indo-European}, sister to Greek, Latin, and the Germanic languages including English. This curriculum treats it, with Latin, as a \textbf{taproot}: its words reach east into South Asia and west into Europe, and the lessons trace both directions.
